\section{From interior statistics to the full law}
\label{pd:section}

We now pass from Theorem~\ref{thm:interior} to ordinary prime averages
of every finite collection of ranked factors.  The argument also shows
why no additional assertion about very small factors or the boundary
of a simplex is needed.
The probability argument uses the size-biased viewpoint of
Donnelly and Grimmett~\cite[Section~2]{DonnellyGrimmett} and the
factorial-measure characterization of
Arratia, Kochman and Miller~\cite[Lemma~2 and Section~3.3]{ArratiaKochmanMiller2014}.
We give the passage, including boundary control and sorting, in full.

\subsection{Factorial measures in the open simplex}

Choose a prime $p$ uniformly from $x<p\leq 2x$, and regard the prime
factors of $p-1$, with their multiplicities, as distinct labelled balls.
A ball corresponding to a prime $q$ has mass
\[
 t=\frac{\log q}{\log(p-1)}.
\]
The masses of all the balls sum to one.  For $d\geq1$, let $\mu_{x,d}$
be the expected counting measure of ordered $d$-tuples of distinct
balls, mapped to their masses.  Distinct balls may correspond to the
same numerical prime.  Set
\begin{equation}\label{pd:open-simplex}
 \mathcal D_d=\left\{(t_1,\ldots,t_d):t_i>0,\quad
                                      \sum_{i=1}^d t_i<1\right\}.
\end{equation}
We first prove the local convergence
\begin{equation}\label{pd:factorial-convergence}
 \int_{\mathcal D_d}\varphi\,\dd\mu_{x,d}
 \longrightarrow
 \int_{\mathcal D_d}\varphi(t)\prod_{i=1}^d\frac{\dd t_i}{t_i}
 \qquad\bigl(\varphi\in C_c(\mathcal D_d)\bigr).
\end{equation}

Start with a rectangle
\begin{equation}\label{pd:admissible-rectangle}
 R=\prod_{i=1}^d(a_i,b_i],\qquad
 0<a_i<b_i,\qquad\sum_i b_i<1.
\end{equation}
Use in Theorem~\ref{thm:interior} the prime slots
$x^{a_i}<\ell_i\leq x^{b_i}$.  They satisfy its hypotheses for a fixed
positive $\eps$.  The reciprocal sum over distinct numerical primes
obeys
\begin{equation}\label{pd:rectangle-main-term}
 C_x=\sum_{\substack{\ell_i\text{ in their slots}\\
                       \ell_i\text{ distinct}}}\frac1{\ell_1\cdots\ell_d}
 \longrightarrow\prod_{i=1}^d\log\frac{b_i}{a_i}.
\end{equation}
Indeed, Lemma~\ref{lem:prime-estimates} gives the limit in each slot.
Terms with a coincidence in two slots have total reciprocal mass
$O(\sum_{q\geq x^{\min a_i}}q^{-2})=o(1)$, times bounded reciprocal
sums from the other slots.

To remove the smooth cutoff in Theorem~\ref{thm:interior}, approximate
the indicator of the prime dyad from above and below by fixed smooth
functions of $(p-1)/x$.  The tuple count is bounded by a constant
depending on $d$ and $\min a_i$, because an integer of size $O(x)$ has
at most $O(1/\min a_i)$ prime factors exceeding $x^{\min a_i}$.
The prime number theorem bounds the contribution of endpoint strips
of relative width $\eta$ by $O(\eta x/\log x)+o(x/\log x)$.
First let real $x$ tend to infinity with the cutoffs fixed, and then
let $\eta$ decrease to zero.  Since
$\pi(2x)-\pi(x)\sim x/\log x$, this proves the rectangle formula
with masses initially normalized by $\log x$.

The required denominator is $\log(p-1)$.  Uniformly on the dyad,
\begin{equation}\label{pd:denominator}
 \frac{\log(p-1)}{\log x}=1+O\left(\frac1{\log x}\right).
\end{equation}
For fixed sufficiently small $\eta>0$, a rectangle with endpoints
$a_i+\eta,b_i-\eta$ on the $\log x$ scale is therefore contained in
the test in \eqref{pd:admissible-rectangle} on the exact scale;
the latter is contained in the rectangle with endpoints
$a_i-\eta,b_i+\eta$.  Choose $\eta$ small enough that all endpoints
stay positive and the enlarged upper endpoints still sum to less
than one.  Apply the preceding rectangle limits, then let $\eta$
decrease to zero.  This proves the same formula with the exact
denominator, still for distinct numerical primes.

The distinction between numerical primes and labelled balls is
negligible locally.  If every tested coordinate is at least $a>0$,
a difference requires $q^2\mid p-1$ for some $q\geq x^{a/2}$, for
all sufficiently large $x$.  The number of possible integers $p-1$
in the dyad is at most
\begin{equation}\label{pd:square-exceptions}
 \sum_{q\geq x^{a/2}}\left\lfloor\frac{2x}{q^2}\right\rfloor
 \ll x^{1-a/2}.
\end{equation}
Here and below a sum over primes may be bounded by the corresponding
sum over integers.  Each such integer contributes at most a constant
number of tested tuples: there are at most $1/a$ balls of mass at
least $a$.  Division by $\pi(2x)-\pi(x)$ makes
\eqref{pd:square-exceptions} negligible.  We have proved
\begin{equation}\label{pd:rectangle-limit}
 \mu_{x,d}(R)\longrightarrow
 \prod_i\log(b_i/a_i)
 =\int_R\prod_i\frac{\dd t_i}{t_i}.
\end{equation}

For completeness, these restricted rectangles determine all the
local limits claimed in \eqref{pd:factorial-convergence}.  A compact
set in $\mathcal D_d$ has all coordinates at least some $a>0$ and
its coordinate sum at most $1-a$, after decreasing $a$ if necessary.
Its $\mu_{x,d}$-mass is bounded uniformly by $a^{-d}$, since the
underlying masses sum to one.  A sufficiently fine rectangular grid
on a slightly larger compact set consists of boxes with positive
lower endpoints and upper endpoints summing to less than one.
Apply \eqref{pd:rectangle-limit} to these finitely many boxes.
Upper and lower step approximations to a continuous function have
an error bounded by its modulus of continuity times the uniformly
bounded local mass.  Refining the fixed grid after taking the
$x$-limit proves \eqref{pd:factorial-convergence}.  This argument
uses no bound for the total factorial mass near zero.

\subsection{Size-biased sampling and the boundary}

Given the balls, sample them successively without replacement,
choosing at each step a remaining ball with probability equal to its
mass divided by the total remaining mass.  Write $T_{x,i}$ for the
mass selected at step $i$, and set all subsequent values to zero
after the finite collection is exhausted.  Let $\nu_{x,d}$ be the
law of $(T_{x,1},\ldots,T_{x,d})$ on $[0,1]^d$.
For an ordered tuple lying in $\mathcal D_d$, its conditional
probability of being the first $d$ draws is
\begin{equation}\label{pd:size-bias-weight}
 w_d(t)=\prod_{i=1}^d
             \frac{t_i}{1-t_1-\cdots-t_{i-1}}.
\end{equation}
Consequently, for $\varphi\in C_c(\mathcal D_d)$,
\[
 \int\varphi\,\dd\nu_{x,d}
 =\int\varphi w_d\,\dd\mu_{x,d}.
\]
The function $w_d$ is continuous and bounded on every compact subset
of $\mathcal D_d$, so \eqref{pd:factorial-convergence} gives the
limiting interior density
\begin{equation}\label{pd:size-bias-density}
 h_d(t)=\prod_{i=1}^d
                   (1-t_1-\cdots-t_{i-1})^{-1}.
\end{equation}

This density has total mass one, as can be checked directly.  Put
\begin{equation}\label{pd:triangular-map}
 A_i=\frac{t_i}{1-t_1-\cdots-t_{i-1}},\qquad
 t_i=A_i\prod_{r<i}(1-A_r).
\end{equation}
These formulas give inverse bijections between $\mathcal D_d$ and
$(0,1)^d$.  The inverse map is triangular and its Jacobian is
\[
 \prod_{i=1}^d\prod_{r<i}(1-A_r)
 =\prod_{i=1}^d(1-t_1-\cdots-t_{i-1}).
\]
It follows that $h_d(t)\,\dd t=\dd A_1\cdots\dd A_d$.
In particular,
\begin{equation}\label{pd:total-mass}
 \int_{\mathcal D_d}h_d(t)\,\dd t=1.
\end{equation}
Taking $U_i=1-A_i$ identifies this measure with the first $d$ stick
fragments $B_i=(\prod_{r<i}U_r)(1-U_i)$ in Theorem~\ref{thm:main},
where the $U_i$ are independent uniform random variables.

Equation~\eqref{pd:total-mass} also excludes any escaped boundary
mass.  Given $\eta>0$, choose $0\leq\chi\leq1$ in
$C_c(\mathcal D_d)$ with $\int\chi h_d>1-\eta$.
Interior convergence yields
$\int\chi\,\dd\nu_{x,d}>1-2\eta$ for large $x$.  Thus at most
$2\eta$ of the probability lies outside $\supp\chi$.  This includes
all boundary faces, exhaustion events, and configurations with
padded zeros.  For any continuous $H$ on $[0,1]^d$, apply interior
convergence to $\chi H$; the integrals of $(1-\chi)H$ under the
two probability measures have total absolute value at most
$3\eta\norm{H}_\infty$ in the limit superior.  Letting $\eta$
decrease to zero proves
\begin{equation}\label{pd:draw-convergence}
 (T_{x,1},\ldots,T_{x,d})
 \ \Longrightarrow\ (B_1,\ldots,B_d)
 \quad\text{on }[0,1]^d.
\end{equation}

\subsection{Sorting and ordinary prime averages}

Let
\[
 R_{x,d}=1-\sum_{i=1}^dT_{x,i},\qquad
 R_d=1-\sum_{i=1}^dB_i=\prod_{i=1}^dU_i.
\]
By \eqref{pd:draw-convergence},
\begin{equation}\label{pd:unseen-mass}
 \E R_{x,d}\longrightarrow\E R_d=2^{-d}.
\end{equation}
The decreasing sequence $R_d$ has a limit whose expectation is
zero, by monotone convergence applied to $1-R_d$.  Hence
$\sum_{i\geq1}B_i=1$ almost surely.  Its decreasing rearrangement
is therefore well-defined and has total mass one.

For a finite or summable nonnegative collection with total mass one,
let $V_i$ be its decreasing rearrangement.  Sort any selected $d$
members and pad the resulting list with zeros, writing $S_{d,i}$.
If the unselected mass is $r$, then
\begin{equation}\label{pd:sorting-bound}
 0\leq V_i-S_{d,i}\leq r\qquad(i\geq1).
\end{equation}
The first inequality follows because removing members cannot
increase any order statistic.  To see the second, if $V_i>r$,
every member at least $V_i$ must have been selected, since each
unselected member is at most $r$.  Thus $S_{d,i}=V_i$ in that
case.  If $V_i\leq r$, the asserted bound is immediate.

Fix the target dimension $k$ and a bounded continuous function
$F:[0,1]^k\to\R$.  Let $\omega_F(\eta)$ be its modulus of
continuity for the maximum norm.  Applying \eqref{pd:sorting-bound}
and Markov's inequality gives
\begin{equation}\label{pd:sorting-test-bound}
 \left|\E F(V_1,\ldots,V_k)
       -\E F(S_{d,1},\ldots,S_{d,k})\right|
 \leq\omega_F(\eta)
       +2\norm{F}_\infty\frac{\E r}{\eta}.
\end{equation}
For fixed $d$, sorting $d$ coordinates and padding is a continuous
map; for example, each order statistic is a finite maximum of
finite minima.  Equation~\eqref{pd:draw-convergence} therefore gives
convergence of the finite sorting test.  Apply
\eqref{pd:sorting-test-bound} both to the factor balls and to the
stick fragments.  First let real $x$ tend to infinity, use
\eqref{pd:unseen-mass}, then let $d$ tend to infinity for fixed
$\eta$, and finally let $\eta$ decrease to zero.  We obtain
\begin{equation}\label{pd:dyad-law}
 \frac1{\pi(2x)-\pi(x)}
 \sum_{x<p\leq2x}F(V_1(p),\ldots,V_k(p))
 \longrightarrow\E F(L_1,\ldots,L_k).
\end{equation}
Every step used limits through all real $x$.

\begin{proof}[Completion of the proof of Theorem~\ref{thm:main}]
For a real upper endpoint $X$, fix an integer $J_0\geq1$ and
partition $(X/2^{J_0},X]$ into the $J_0$ dyads
$(X/2^j,X/2^{j-1}]$.  Each dyad scale tends to infinity with $X$,
so \eqref{pd:dyad-law}, applied a finite number of times, shows
that the weighted sum over these dyads has the asserted limiting
average.  The discarded primes below $X/2^{J_0}$ have proportion
\[
 \frac{\pi(X/2^{J_0})}{\pi(X)-1}=2^{-J_0}+o(1)
\]
by the prime number theorem.  Their effect on the discrepancy from
the limiting expectation is at most
$2\norm{F}_\infty(2^{-J_0}+o(1))$.
Take $X\to\infty$, then $J_0\to\infty$.  Removing the prime $2$
changes the normalization and sum by $o(1)$.  This proves the
statement for ordinary equal weighting of $3\leq p\leq X$ and
for every real upper endpoint tending to infinity.
\end{proof}
